\chapter{演習問題}

\newcommand{\schoolCalender}[1]{%
\ifcase#1%
\or 4月14日%
\or 4月21日%
\or 4月28日%
\or 5月12日%
\or 5月14日(6)%
\or 5月19日%
\or 5月26日%
\or 6月2日%
\or 6月9日%
\or 6月16日%
\or 6月23日%
\or 6月30日%
\or 7月7日%
\or 7月14日%
\or 7月28日%
\or 7月31日%
\else %
\fi
}



%% \section*{予習の進め方について}
%% この講義は, 所謂反転授業の形式で行う.
%% Moodleにある動画資料を予習し講義に臨み,
%% 講義時間には簡単な解説および問題演習を行う.

%% 予習用のノートと講義での解説のノートは分けたほうが良いと思われる.

%% 予習はできるときに自分のペースで先に進むとよいと思う.
%% 遅くとも, 以下のように予習を済ませることが望ましい:
%% \begin{enumerate}
%% \item \schoolCalender{2}の講義までに,
%%   1. 行列に関する用語の定義
%%   まで.
%% \item \schoolCalender{3}の講義までに,
%%   2. 行列の演算
%%   まで.
%% \item \schoolCalender{4}の講義までに,
%%   3. 行列式と逆行列
%%   まで.
%% %%%%%%%%%% この年は, 5週目がイレギュラーな曜日
%% \item \schoolCalender{6}の講義までに,
%%   4. 連立方程式 (1)
%%   まで.
%% \item \schoolCalender{7}の講義までに,
%%   5. 連立方程式 (2)
%%   まで.
%% \item \schoolCalender{9}の講義までに,
%%   6. 一次独立性と基底
%%   まで.
%% \item \schoolCalender{10}の講義までに,
%%   7. 内積とノルム
%%   まで.
%% \item \schoolCalender{11}の講義までに,
%%   8. 平面上の線形変換
%%   まで.
%% \item \schoolCalender{13}の講義までに,
%%   9. 固有値と固有ベクトル
%%   まで.
%% \item \schoolCalender{14}の講義までに,
%%   10. 固有値と固有ベクトル(2)
%%   まで.
%% \end{enumerate}

%% \newpage

\section{1/15回目: \schoolCalender{1}}
\Cref{sec:def:mat}に関する内容.
% 行列に関する用語 (型, 成分, 行, 列, 転置) を対面で解説する.
\subsection{問題}
\begin{quiz}
  次の行列$A$に対し, $\transposed{A}$を求めよ:
  \begin{enumerate}
  \item $A=\begin{pmatrix}
    1&2\\
    3&4\\
    5&6
    \end{pmatrix}$.
  \item $A=\begin{pmatrix}
    1\\
    2
    \end{pmatrix}$.
  \end{enumerate}
\end{quiz}

\subsection{WebWork}
\begin{enumerate}
\item Set01 Problem 1 (行列の用語の確認) を取り組め.
\item Set01 Problem 3 (行列の成分) を取り組め.
\item Set01 Problem 7 (行列の相等) を取り組め.
\end{enumerate}

%\endinput
\newpage
\section{2/15回目: \schoolCalender{2}}
\Cref{sec:def:matop}に関する内容.
% 零行列, 単位行列, クロネッカーのデルタについて対面で解説する.
% 対面で和とスカラー倍を定義する.
% 予習は   1. 行列に関する用語の定義
\subsection{問題}
\begin{quiz}
  次の$A$, $B$に対し, $A+B$を計算せよ:
  \begin{enumerate}
  \item
    $A=\begin{pmatrix}1&9\\2&-3\end{pmatrix}$,
    $B=\begin{pmatrix}4&0\\5&7\end{pmatrix}$.
  \item 
    $A=\begin{pmatrix}2\\\pi\end{pmatrix}$,
    $B=\begin{pmatrix}e\\5\end{pmatrix}$.
  \end{enumerate}
\end{quiz}
\begin{quiz}
  次の$A$, $c$に対し, $cA$を計算せよ:
  \begin{enumerate}
  \item
    $c=-4$,
    $A=\begin{pmatrix}1&\sqrt{-1}\\2&-3\end{pmatrix}$.
  \item
    $c=\pi$,
    $A=\begin{pmatrix}e^2\\3\end{pmatrix}$.
  \end{enumerate}
\end{quiz}
\subsection{WebWork}
\begin{enumerate}
\item Set01 Problem 2 (行列の用語の確認) を取り組め.
\item Set01 Problem 4--6 (クロネッカーのデルタ) を取り組め.
\item Set01 Problem 8 (成分の比較) を取り組め.
\item Set01 Problem 14 (特別な形の行列) を取り組め.
\item Set02 Problem 2--3, 11--12, 20, 23--24 (和)を取り組め.
\item Set02 Problem 33--37 (和, スカラー倍)を取り組め.
\end{enumerate}


%\endinput
\newpage
\section{3/15回目: \schoolCalender{3}}
\Cref{chap:mat}に関する内容.
% 対面で行列の積について定義する.
% 予習は 2. 行列の演算
\subsection{問題}
\begin{quiz}
  以下の$A$, $B$に対し, $AB$を求めよ:
  \begin{enumerate}
  \item
    $A=\begin{pmatrix}5&6\\7&8\end{pmatrix}$,
    $B=\begin{pmatrix}e&\pi\\\sqrt{2}&\sqrt{3}\end{pmatrix}$.
  \item
    $A=\begin{pmatrix}1&2&0\\2&3&4\end{pmatrix}$,
    $B=\begin{pmatrix}5&6\\7&8\\-1&5\end{pmatrix}$.
  \item
    $A=\begin{pmatrix}2&2&3\\1&3&4\end{pmatrix}$,
    $B=\begin{pmatrix}5\\7\\6\end{pmatrix}$.
  \item
    $A=\begin{pmatrix}2&3\end{pmatrix}$,
    $B=\begin{pmatrix}e\\\pi\end{pmatrix}$.
  \item
    $A=\begin{pmatrix}2\\3\end{pmatrix}$,
    $B=\begin{pmatrix}e&\pi\end{pmatrix}$.
  \end{enumerate}
\end{quiz}
\subsection{WebWork}
\begin{enumerate}
\item Set02 Problem 1(和と積が定義できるための条件)を取り組め.
\item Set02 Problem 4--8, 13--14, 21--22, 25--26 (積)を取り組め.
\item Set02 Problem 9--10, 15--16, 27--28 (行列の和と積) を取り組め.
\item Set01 Problem 9 (転置) を取り組め.
\item Set01 Problem 10--13 (対称行列, 歪対称行列) を取り組め.
\end{enumerate}

%\endinput
\newpage
\section{4/15回目: \schoolCalender{4}}
\Cref{chap:inverse}に関する問題.
%  3. 行列式と逆行列
\subsection{問題}
\begin{quiz}
  次の行列$A$が正則かどうか調べ, 正則なら$A$の逆行列を求めよ:
  \begin{enumerate}
    \item $A=\begin{pmatrix}1&2\\3&4\end{pmatrix}$.
    \item $A=\begin{pmatrix}1&1\\1&1\end{pmatrix}$.
    \item $A=\begin{pmatrix}0&1\\1&0\end{pmatrix}$.
    \item $A=\begin{pmatrix}\cos(\theta)&-\sin(\theta)\\\sin(\theta)&\cos(\theta)\end{pmatrix}$.
    \item $A=\begin{pmatrix}t&1\\1&t\end{pmatrix}$.
  \end{enumerate}
\end{quiz}

\subsection{WebWork}
\begin{enumerate}
\item Set03 Problem 1--3 (行列式) を取り組め.
\item Set03 Problem 4 (行列式による方程式) を取り組め.
\item Set03 Problem 5 (行列式と正則性) を取り組め.
\item Set03 Problem 6--10 (逆行列) を取り組め.

\item Set03 Problem 13--14 (逆行列と連立方程式) を取り組め.
\end{enumerate}

%\endinput
\newpage
\section{5/15回目: \schoolCalender{5}}
\Cref{def:det:generic,thm:fund:inverse}に関する問題.
(イレギュラーな曜日に割り振られた回)
\subsection{問題}
\begin{quiz}
  次の行列$A$の行列式を求め,
  正則かどうか判定せよ.
  \begin{enumerate}
  \item
    $A=
    \begin{pmatrix}
      1&2&0&0\\
      0&4&6&9\\
      0&5&6&9\\
      0&6&6&9
    \end{pmatrix}$
  \item
    $A=
    \begin{pmatrix}
      1&0&3&4\\
      0&4&6&9\\
      0&3&5&8\\
      0&1&1&1
    \end{pmatrix}$
\end{enumerate}
\end{quiz}

\begin{quiz}
  次の行列が
  正則かどうか判定し,
  正則なら逆行列を求めよ.
  \begin{enumerate}
  \item
    $A=
    \begin{pmatrix}
      1&2&3\\
      0&4&8\\
      0&4&6
    \end{pmatrix}$
  \item
    $A=
    \begin{pmatrix}
      1&0&3\\
      0&4&4\\
      1&4&7
    \end{pmatrix}$
\end{enumerate}
\end{quiz}


\subsection{WebWork}
\begin{enumerate}
\item Set 05 Problem 1--8 (3次行列式)
\item Set 05 Problem 9--13 (3次逆行列)
\item Set 05 Problem 15--16 (3次逆行列と連立方程式)
\item Set 05 Problem 14 (変数の入った3次逆行列)
\end{enumerate}

%\endinput
\newpage
\section{6/15回目: \schoolCalender{6}}
\Cref{chap:systemoflineq}に関する問題.
%  4 (1). 連立方程式 (1)
% 被約階段行列と行基本変形と階数.
\subsection{問題}
\begin{quiz}
  次の行列$A$に行基本変形をして得られる被役行階段行列を求めよ.
  $A$の階数を答えよ. $A$が正方行列なら, $A$が正則かどうか答えよ.
  \begin{enumerate}
  \item $A=\begin{pmatrix}1&2&3\\1&3&3\\1&2&5\end{pmatrix}$.
  \item $A=\begin{pmatrix}1&2\\1&3\\1&2\end{pmatrix}$.
  \item $A=\begin{pmatrix}1&0&4\\1&1&5\\0&1&2\end{pmatrix}$.
  \item $A=\begin{pmatrix}1&0&4\\0&2&2\\1&2&6\end{pmatrix}$.
  \item $A=\begin{pmatrix}1&0&4\\1&&4\\0&0&1\end{pmatrix}$.
  \end{enumerate}
\end{quiz}


\subsection{WebWork}
\begin{enumerate}
\item Set04 Problem 1--4, 12 (被役階段行列, ランク) を取り組め.
\item Set04 Problem 5 (連立一次方程式の用語) を取り組め.
\item Set06 Problem 1--4, 12 (被役階段行列の定義) を取り組め.
\item Set06 Problem 5--10 (被役階段行列) を取り組め.
\item Set06 Problem 11--14 (被役階段行列, ランク) を取り組め.

\end{enumerate}

%\endinput
\newpage
\section{7/15回目: \schoolCalender{7}}
\Cref{chap:systemoflineq}に関する問題.
%   4 (2). 連立方程式 (2)
% 連立方程式全部
% 連立方程式は行基本変形で被約階段行列に変形すれば解けること.
% 正則性と解の自由度の話もする.
\subsection{問題}
\begin{quiz}
  次の行列$A$を拡大係数行列とする連立一次方程式の解の空間を求めよ:
  \begin{enumerate}
  \item $A=\begin{pmatrix}0&1&0&2&3&0&1&0&5\\0&0&1&2&0&0&0&0&0\\0&0&0&0&0&1&0&0&9\\0&0&0&0&0&0&0&0&0
  \end{pmatrix}$.
  \item $A=\begin{pmatrix}1&2&3\\0&0&0\end{pmatrix}$.
    \item $A=\begin{pmatrix}1&2&3\\2&3&5\end{pmatrix}$.
    \item $A=\begin{pmatrix}1&2&3\\2&4&5\end{pmatrix}$.
    \item $A=\begin{pmatrix}1&2&3\\2&4&6\end{pmatrix}$.
    \item $A=\begin{pmatrix}0&0&0\\0&0&0\end{pmatrix}$.
  \end{enumerate}
\end{quiz}

\begin{quiz}
  次の行列$A$と$\bbb$に対し,
  $\xx$に関する方程式$A\xx=\bbb$の解の空間を求めよ.
  また, 解を持つ場合には解の自由度も答えよ.
  \begin{enumerate}
    \item $A=\begin{pmatrix}1&2\\2&3\end{pmatrix}$,
     $\bbb=\begin{pmatrix}3\\5\end{pmatrix}$.
    \item $A=\begin{pmatrix}1&2\\2&4\end{pmatrix}$,
    $\bbb=\begin{pmatrix}3\\6\end{pmatrix}$.
    \item $A=\begin{pmatrix}0&0\\0&0\end{pmatrix}$,
      $\bbb=\begin{pmatrix}0\\0\end{pmatrix}$.
  \end{enumerate}
\end{quiz}
\subsection{WebWork}
\begin{enumerate}
\item Set04 Problem 6-11 (連立一次方程式) を取り組め.
\end{enumerate}

%\endinput
\newpage
\section{8/15回目: \schoolCalender{8}}
\Cref{chap:systemoflineq}に関する問題.
% 連立一次方程式/続き.
\subsection{問題}
\begin{quiz}
  次の行列$A$を拡大係数行列とする連立一次方程式の解の空間を求めよ:
  \begin{enumerate}
  \item $A=\begin{pmatrix}0&1&0&2&3&0&1&0&5\\0&0&1&2&0&0&0&0&0\\0&0&0&0&0&1&0&0&9\\0&0&0&0&0&0&0&0&0\end{pmatrix}$.
  \end{enumerate}
\end{quiz}
\begin{quiz}
  次の行列$A$に対し$A\xx=\zzero$が$\xx=\zzero$以外の解を持つか判定せよ:
  \begin{enumerate}
  \item $A=\begin{pmatrix}2&1\\1&2\end{pmatrix}$.
  \item $A=\begin{pmatrix}2&1\\2&1\end{pmatrix}$.
  \item $A=\begin{pmatrix}3&1\\2&1\end{pmatrix}$.
  \end{enumerate}
\end{quiz}

\begin{quiz}
  次の行列$A$に対し$A\xx=\zzero$が$\xx=\zzero$以外の解を持つような$c$を求めよ:
  \begin{enumerate}
  \item $A=\begin{pmatrix}2-c&1\\1&2-c\end{pmatrix}$.
  \item $A=\begin{pmatrix}3-c&1\\2&1-c\end{pmatrix}$.
  \end{enumerate}
\end{quiz}

\subsection{WebWork}
\begin{enumerate}
\item Set06 Problem 16--22 (連立一次方程式) を取り組め.
\item Set06 Problem 23 (変数の入った連立一次方程式) を取り組め.
\item Set02 Problem 17--19 (羃の計算).
\item Set02 Problem 29--31 (羃の計算).
\end{enumerate}
以下はすでに演習で扱った内容であるので,
まだ取り組んでいない場合は適宜取り組め.
\begin{enumerate}
\item Set01 Problem 1--14
\item Set02 Problem 1--16, 20--28  33-37
\item Set03 Problem 1--10, 13--14
\item Set04 Problem 1--12
\item Set06 Problem 1--14 16-22
\end{enumerate}


%\endinput
\newpage
\section{9/15回目: \schoolCalender{9}}
\Cref{sec:vect}に関する問題.
% ノルム, 内積は対面でやる.
% 一次独立性と基底
\subsection{問題}
\begin{quiz}
  次のベクトルの組が一次独立か判定せよ:
  \begin{enumerate}
    \item $(\begin{pmatrix}1\\2\\3\end{pmatrix},\begin{pmatrix}1\\2\\4\end{pmatrix})$.
    \item $(\begin{pmatrix}1\\1\end{pmatrix},\begin{pmatrix}2\\2\end{pmatrix},\begin{pmatrix}3\\4\end{pmatrix})$.
    \item $(\begin{pmatrix}1\\1\end{pmatrix},\begin{pmatrix}2\\5\end{pmatrix})$.
    \item $(\begin{pmatrix}1\\1\end{pmatrix})$.
    \item $(\begin{pmatrix}0\\0\end{pmatrix})$.
  \end{enumerate}
\end{quiz}


\begin{quiz}
  次のベクトルの組$\RR^2$の基底であるか判定せよ:
  \begin{enumerate}
    \item $(\begin{pmatrix}1\\2\end{pmatrix},\begin{pmatrix}1\\3\end{pmatrix})$.
    \item $(\begin{pmatrix}1\\2\end{pmatrix},\begin{pmatrix}2\\4\end{pmatrix})$.
    \item $(\begin{pmatrix}1\\2\end{pmatrix})$.
    \item $(\begin{pmatrix}1\\0\end{pmatrix},\begin{pmatrix}0\\1\end{pmatrix},\begin{pmatrix}1\\1\end{pmatrix})$.
  \end{enumerate}
\end{quiz}


\subsection{WebWork}
\begin{enumerate}
  \item Set07 Problem 1--3 (一次独立性, 基底) を取り組め.
  \item Set07 Problem 4--8 (内積とノルム) を取り組め.
\end{enumerate}

\endinput
\newpage
\section{10/15回目: \schoolCalender{10}}
\subsection{問題}
%   7. 内積とノルム
\begin{quiz}
\end{quiz}
\subsection{WebWork}
\begin{enumerate}
  \item Set Problem を取り組め.
\end{enumerate}

\endinput
\newpage
\section{11/15回目: \schoolCalender{11}}
\subsection{問題}
\begin{quiz}
\end{quiz}

\subsection{WebWork}
%   8. 平面上の線形変換
\begin{enumerate}
  \item Set Problem を取り組め.
\end{enumerate}

\endinput
\newpage
\section{12/15回目: \schoolCalender{12}}
\subsection{問題}

\subsection{WebWork}
\begin{enumerate}
  \item Set Problem を取り組め.
\end{enumerate}

\endinput
\newpage
\section{13/15回目: \schoolCalender{13}}
\subsection{問題}

\subsection{WebWork}
\begin{enumerate}
  \item Set Problem を取り組め.
\end{enumerate}


\endinput
\newpage
\section{14/15回目: \schoolCalender{14}}
\subsection{問題}
\begin{quiz}
\end{quiz}

\subsection{WebWork}
\begin{enumerate}
  \item Set Problem を取り組め.
\end{enumerate}

\endinput
\newpage
\section{15/15回目: \schoolCalender{15}}
\subsection{問題}
\begin{quiz}
\end{quiz}

\subsection{WebWork}
\begin{enumerate}
  \item Set Problem を取り組め.
\end{enumerate}


\begin{enumerate}
%
\item Set02 Problem 17--19, 29, 30--31 (行列の冪) を取り組め.
\item Set02 Problem 32 (ケーリーハミルトンの定理) を取り組め.
\item Set03 Problem 11--12 (行列式と面積) を取り組め.
\item Set06 Problem 24 (クラメルの公式) を取り組め.
\end{enumerate}

\endinput
\newpage
